\documentclass[a4paper,12pt,twocolumn]{article}

\usepackage[landscape,top=1cm,bottom=1.4cm,left=1cm,right=1cm,columnsep=1cm]{geometry}
\usepackage[fleqn]{amsmath}
\usepackage{amssymb}
\usepackage{fontspec}
\usepackage{enumitem}
\usepackage{xeCJK}
\usepackage[hidelinks]{hyperref}

\setlength{\mathindent}{1.5em}

% Set fonts
\setmainfont{Times New Roman}
\setsansfont{Arial}
\setmonofont{JetBrains Mono}

% Chinese font support
\setCJKmainfont{Iansui}[
  BoldFont=Iansui,
  AutoFakeBold=3.17
]

\pagestyle{plain}
\setlength{\columnseprule}{0.2pt}

\newcommand{\examdate}{2026/04/17}
\newlength{\problemnumberwidth}
\settowidth{\problemnumberwidth}{\textbf{69.}\ }
\newcommand{\problemtitle}[2]{%
  \hypertarget{#1}{}%
  \par\vspace{0.4em}%
  \noindent{\large\bfseries #2}\par
  \vspace{0.3em}%
  \hrule
  \vspace{0.9em}%
}
\newcommand{\worksheettoc}{%
  \noindent{\Large\bfseries Contents}\par
  \vspace{0.3em}%
  \hrule
  \vspace{0.9em}%
  \begin{tabular}{@{}p{0.26\columnwidth}p{0.66\columnwidth}@{}}
    \hyperlink{sec:13-1}{\textbf{13-1}} & \hyperlink{sec:13-1}{41, 42, 43, 44, 45, 46, 57, 59, 60} \\
    \hyperlink{sec:13-2}{\textbf{13-2}} & \hyperlink{sec:13-2}{24, 28, 29, 35, 52} \\
  \end{tabular}\par
}
\newcommand{\problem}[2]{%
  \par\noindent
  \hangindent=\problemnumberwidth
  \hangafter=1
  \makebox[\problemnumberwidth][l]{\textbf{#1.}}#2\par
  \vspace{1.4em}%
}
\newlist{subparts}{enumerate}{1}
\setlist[subparts]{%
  label = (\alph*),
  labelindent=\problemnumberwidth,
  leftmargin=*,
  labelsep=0.5em,
  align=left,
  itemsep=0.2em,
  topsep=0.3em,
  parsep=0pt,
  partopsep=0pt
}

\begin{document}

\twocolumn[
{\centering
\Large\bfseries Calculus Problems\par
\vspace{0.5em}
\noindent\hfill\normalsize\normalfont \examdate\par
\vspace{0.8em}
}
]

\worksheettoc
\vfill\null
\newpage

\problemtitle{sec:13-1}{13-1}

\problem{41}{Graph the curve with the given vector equation. Make sure you choose a parameter domain and viewpoints that reveal the true nature of the curve.
\[
r(t)=\langle \cos t\sin 2t,\ \sin t\sin 2t,\ \cos 2t\rangle
\]}

\problem{42}{Graph the curve with the given vector equation. Make sure you choose a parameter domain and viewpoints that reveal the true nature of the curve.
\[
r(t)=\langle te^t,\ e^{-t},\ t\rangle
\]}

\problem{43}{Graph the curve with the given vector equation. Make sure you choose a parameter domain and viewpoints that reveal the true nature of the curve.
\[
r(t)=\left\langle \sin 3t\cos t,\ \frac{1}{4}t,\ \sin 3t\sin t\right\rangle
\]}

\problem{44}{Graph the curve with the given vector equation. Make sure you choose a parameter domain and viewpoints that reveal the true nature of the curve.
\[
r(t)=\langle \cos(8\cos t)\sin t,\ \sin(8\cos t)\sin t,\ \cos t\rangle
\]}

\problem{45}{Graph the curve with the given vector equation. Make sure you choose a parameter domain and viewpoints that reveal the true nature of the curve.
\[
r(t)=\langle \cos 2t,\ \cos 3t,\ \cos 4t\rangle
\]}

\problem{46}{Graph the curve with parametric equations
\[
x=\sin t \qquad y=\sin 2t \qquad z=\cos 4t
\]
Explain its shape by graphing its projections onto the three coordinate planes.}

\problem{57}{Intersection and Collision. The trajectories of two particles are given by the vector functions
\[
r_1(t)=\langle t^2,\ 7t-12,\ t^2\rangle \qquad r_2(t)=\langle 4t-3,\ t^2,\ 5t-6\rangle
\]
for \(t\ge 0\). Do the particles collide?}

\problem{59}{%
\begin{subparts}
\item Graph the curve with parametric equations
\[
x=\frac{27}{26}\sin 8t-\frac{8}{39}\sin 18t
\]
\[
y=-\frac{27}{26}\cos 8t+\frac{8}{39}\cos 18t
\]
\[
z=\frac{144}{65}\sin 5t
\]
\item Show that the curve lies on the hyperboloid of one sheet \(144x^2+144y^2-25z^2=100\).
\end{subparts}}

\problem{60}{Trefoil Knot. The view of the trefoil knot shown in Figure 9 is accurate, but it doesn't reveal the whole story. Use the parametric equations
\[
x=(2+\cos 1.5t)\cos t
\]
\[
y=(2+\cos 1.5t)\sin t
\]
\[
z=\sin 1.5t
\]
to sketch the curve by hand as viewed from above, with gaps indicating where the curve passes over itself. Start by showing that the projection of the curve onto the \(xy\)-plane has polar coordinates \(r=2+\cos 1.5t\) and \(\theta=t\), so \(r\) varies between \(1\) and \(3\). Then show that \(z\) has maximum and minimum values when the projection is halfway between \(r=1\) and \(r=3\).
{\medskip\par
\hangindent=\problemnumberwidth
\hangafter=1
\noindent\hspace*{\problemnumberwidth}When you have finished your sketch, use a computer to draw the curve with viewpoint directly above and compare with your sketch. Then plot the curve from several other viewpoints. You can get a better impression of the curve if you plot a tube with radius \(0.2\) around the curve. (Use the \texttt{tubeplot} command in Maple or the \texttt{tubecurve} or \texttt{Tube} command in Mathematica.)}}

\problemtitle{sec:13-2}{13-2}

\problem{24}{If \(r(t)=\langle e^{2t},\ e^{-3t},\ t\rangle\), find \(r'(0)\), \(T(0)\), \(r''(0)\), and \(r'(0)\times r''(0)\).}

\problem{28}{Find parametric equations for the tangent line to the curve with the given parametric equations at the specified point.
\[
x=\sqrt{t^2+3}, \qquad y=\ln(t^2+3), \qquad z=t; \qquad (2,\ln 4,1)
\]}

\problem{29}{Find a vector equation for the tangent line to the curve of intersection of the cylinders \(x^2+y^2=25\) and \(y^2+z^2=20\) at the point \((3,4,2)\).}

\problem{35}{The curves \(r_1(t)=\langle t,\ t^2,\ t^3\rangle\) and \(r_2(t)=\langle \sin t,\ \sin 2t,\ t\rangle\) intersect at the origin. Find their angle of intersection correct to the nearest degree.}

\problem{52}{If \(r(t)=u(t)\times v(t)\), where \(u(2)=\langle 1,\ 2,\ -1\rangle\), \(u'(2)=\langle 3,\ 0,\ 4\rangle\), and \(v(t)=\langle t,\ t^2,\ t^3\rangle\), find \(r'(2)\).}

\end{document}
